\section{Related work and the limits of transferring evidence}
\label{sec:related-work}

The relevant literature spans several different target labels: synthesized
speech, converted voices, lyrical authorship, source-separation fidelity and
whole-recording music generation. An interpretable descriptor may be
motivated by one task without being validated for another. This thesis
therefore distinguishes the origin of a hypothesis from the external assay
that qualifies its measurement and the held-source experiment that tests
its predictive transfer.

\subsection{Reconstruction artifacts versus generator transfer}
Afchar, Meseguer-Brocal and Hennequin compare original music with matched
autoencoder reconstructions, reducing content and file-encoding confounds
\cite{afchar2025}. Their controlled representation comparison reports
99.8\% accuracy for amplitude and 99.6\% for phase input. These are learned
detector results on that reconstruction task, not validation of the present
hand-designed F statistics. Their subsequent discussion of audio manipulation
and unseen decoders illustrates why near-perfect internal discrimination
is insufficient for deployment. We use this distinction to motivate both
the original spectral/separator controls and the separate source-heldout
evaluation; we do not reproduce their classifier or infer that its scores
are comparable to our macro AUC or balanced accuracy.

\subsection{Musical organization as measurement and representation}
Wu and Yang analyze generated jazz lead sheets using pitch-class, groove,
chord-progression and structural statistics, alongside listening and
continuation tasks \cite{wu2020}. This provides a precedent for examining
musical organization quantitatively. Its symbolic jazz setting does not
establish a universal direction for chroma entropy or repetition in modern
waveform generators. Our H and M implementations are explicitly specified
audio descriptors, not replicas of every measure in that study.

Kim and Go's Segment Transformer uses beat-aware segmentation and learned
inter-segment relationships for music-generation detection
\cite{kim2025}. The work evaluates short and full-audio systems on
FakeMusicCaps and SONICS. It supports considering longer-context information,
but a learned model's improvement does not isolate a hand-designed recurrence
scalar as its explanation. In our experiments, M's standalone and incremental
effects are measured separately, with duration eligibility preserved. The
published model is neither retrained nor added to the current comparison.

\subsection{Breathing evidence is task-specific}
Mostaani and Magimai Doss study embeddings from networks trained to estimate
speech breathing patterns in logical-access presentation attack detection
\cite{mostaani2022}. Their ASVspoof 2019 results distinguish natural from
synthesized speech, but do not show the same distinction for voice conversion.
This is relevant motivation for investigating respiration-related cues, not
validation of event-level breath localization in singing mixtures. The
present project's failed breath-event precision/F1 gate consequently remains
informative even though earlier speech work reports successful discrimination.
An embedding result and a breath-event counter have different measurement
units and failure modes.

\subsection{Synthesized stereo is not native-stereo music evidence}
Liu et al.'s M2S-ADD first converts mono input to stereo using a pretrained
synthesizer and then applies a dual-branch detector \cite{liu2023}. Their
reported evaluation concerns ASVspoof2019. This does not establish that
native stereo in generated songs has an intrinsic spatial inconsistency.
Our SC family instead describes recorded channel relationships and must be
interpreted alongside channel provenance, deliberately created width and
the matched feature-ablation results. A classifier's use of two channels
does not alone identify the physical reason for its prediction.

\subsection{Lyrics authorship and waveform generation}
Frohmann et al.'s \emph{Double Entendre} combines automatically transcribed
sung lyrics with speech representations in a late-fusion system
\cite{frohmann2025}. This provides a complementary approach using lyrical
information, with robustness and generalization motivating its design.
However, whether lyrics were generated and whether a recording was generated
are not identical labels: human-written lyrics can condition synthesized
audio, and generated lyrics can be performed by people. We do not treat
lyrical evidence as ground truth for the present waveform-level labels or
compare the paper's reported performance directly with our source-heldout
scores. No lyrics model is added in this closeout.

\subsection{Position of the present experiments}
The contribution examined here is not an assertion that one additional
artifact universally separates human and generated music. It is the
operationalization and testing of a finite set of interpretable descriptors,
including explicit failures, availability limits and conditional predictive
increments. The same external scientific gate, scalar definition and
source-aware evaluation must be maintained even when a different choice
would improve a displayed score. Literature-derived motivations cannot
override observed failures of those tests.

The records above were rechecked against primary manuscripts or conference
pages during editorial integration. They establish the stated task scopes;
their experiments were not replicated here. The larger retained literature
review and source-verification notes remain supporting provenance, not a
claim that every proposed candidate entered this thesis's experiments.
